\documentclass[11pt]{article}
\usepackage[a4paper,margin=27mm]{geometry}
\usepackage[T1]{fontenc}
\usepackage{lmodern,microtype,amsmath,amssymb,amsthm,mathtools,enumitem}
\usepackage[colorlinks=true,linkcolor=blue,citecolor=blue,urlcolor=blue]{hyperref}

\newtheorem{theorem}{Theorem}[section]
\newtheorem{lemma}[theorem]{Lemma}
\newtheorem{proposition}[theorem]{Proposition}
\newtheorem{corollary}[theorem]{Corollary}
\theoremstyle{definition}
\newtheorem{definition}[theorem]{Definition}
\newenvironment{axiom}[1]{\par\medskip\noindent\textbf{Axiom #1.}\ \itshape}{\par\medskip}

\newcommand{\M}{\mathcal M}
\newcommand{\C}{\mathcal C}
\newcommand{\U}{\mathcal U}
\newcommand{\K}{\mathfrak K}
\newcommand{\Lf}{\mathfrak L}
\newcommand{\one}{\mathbf 1}
\newcommand{\R}{\mathbb R}
\newcommand{\N}{\mathbb N}
\newcommand{\supp}{\operatorname{supp}}
\newcommand{\interior}{\operatorname{int}}
\newcommand{\bd}{\partial}
\newcommand{\cor}{\mathrm{cor}}
\newcommand{\sil}{\mathrm{sil}}
\newcommand{\wro}{\mathrm{wro}}
\newcommand{\con}{\mathrm{con}}
\newcommand{\AV}{\mathrm{AV}}
\newcommand{\WSL}{\mathrm{WSL}}
\newcommand{\CONC}{\mathrm{CONC}}
\newcommand{\GU}{\mathrm{G}}
\setlist{itemsep=1pt,topsep=3pt}

\title{Coherence and Dominance for Sets of Desirable Gambles}
\author{Jason Konek}
\date{}

\begin{document}
\maketitle

\begin{abstract}
[Abstract to be written by the author.]
\end{abstract}

\section{Setting}\label{sec:setting}

\begin{definition}[Gambles]\label{def:gambles}
Let $\Omega$ be a nonempty set. Let $E$ be a closed linear subspace of the space of bounded real functions on $\Omega$, with the norm $\|g\|\coloneqq\sup_{\omega\in\Omega}|g(\omega)|$, such that $E$ is separable and the constant function $\one$ belongs to $E$. Elements of $E$ are \emph{gambles}. For $g,h\in E$ write $g\ge h$ if $g(\omega)\ge h(\omega)$ for every $\omega\in\Omega$. Put
\[
E_+\coloneqq\{g\in E:g\ge0\},\qquad E_{++}\coloneqq\{g\in E:\inf_{\Omega}g>0\},\qquad E_{--}\coloneqq -E_{++}.
\]
Elements of $E_{++}$ are \emph{uniformly positive} and elements of $E_{--}$ are \emph{uniformly negative}. For $c\in E$ and $r>0$, $B(c,r)\coloneqq\{g\in E:\|g-c\|<r\}$. For $A\subseteq E$, $-A\coloneqq\{-g:g\in A\}$; $A$ is \emph{symmetric} if $A=-A$. For a set $R\subseteq E$, $\interior R$ and $\bd R$ denote its interior and boundary.
\end{definition}

\begin{definition}[Judgments and problem classes]\label{def:judgments}
For options $a,b$ of a binary decision problem with $a-b=g\in E$, the gamble $g$ is the \emph{judgment} that $a$ is preferred to $b$, and $-g$ is the judgment that $b$ is preferred to $a$. For $g\in E$, the \emph{problem class} of $g$ is the set $[g]\coloneqq\{g,-g\}$; it indexes all binary problems whose option difference is $g$ or $-g$. A set of problem classes is represented by the symmetric set of gambles that is the union of its members.
\end{definition}

\begin{definition}[Frequency measure]\label{def:frequency}
$\mu$ is a finite Borel measure on $E$ that is \emph{symmetric}: $\mu(-A)=\mu(A)$ for every Borel $A$. For a symmetric measurable $A\subseteq E$, $\pi(A)\coloneqq\tfrac12\mu(A)$ is the frequency of the problem classes contained in $A$. For a measurable $A$ with $A\cap(-A)=\varnothing$, $\mu(A)=\pi(A\cup(-A))$ is the frequency of the problem classes that meet $A$. ``Measurable'' means measurable for the completion of $\mu$, and ``null'' means of $\mu$-measure zero.
\end{definition}

\begin{definition}[Models]\label{def:models}
A \emph{model} is a measurable set $X\subseteq E$, two models being identified when their symmetric difference is null. $\M$ is the set of models. For a model $X$:
\begin{gather*}
U(X)\coloneqq\{g\in X:-g\notin X\},\qquad K(X)\coloneqq\{g\in X:-g\in X\},\\
S(X)\coloneqq\{g\in E:g\notin X,\ -g\notin X\}.
\end{gather*}
$U(X)$ is the set of judgments that $X$ \emph{uniquely recommends}, $K(X)$ the set of judgments on which $X$ is \emph{confused}, and $S(X)$ the set of judgments on which $X$ is \emph{silent}. $S(X)$ and $K(X)$ are symmetric.
\end{definition}

\begin{definition}[Responses at a world]\label{def:responses}
Fix $\omega\in\Omega$ and $g\in E$ with $g(\omega)\ne0$. Let $g^\omega$ be the element of $[g]$ with $g^\omega(\omega)>0$. The \emph{response} of a model $X$ to $[g]$ at $\omega$ is
\[
\rho_\omega(X,g)\coloneqq
\begin{cases}
\cor &\text{if } g^\omega\in X,\ -g^\omega\notin X,\\
\sil &\text{if } g^\omega\notin X,\ -g^\omega\notin X,\\
\wro &\text{if } g^\omega\notin X,\ -g^\omega\in X,\\
\con &\text{if } g^\omega\in X,\ -g^\omega\in X.
\end{cases}
\]
$\rho_\omega(X,g)=\rho_\omega(X,-g)$. Define the total order $\cor\rhd\sil\rhd\wro\rhd\con$ on the four responses, and write $a\unrhd b$ for ($a\rhd b$ or $a=b$).
\end{definition}

\begin{definition}[Changed classes, distance]\label{def:changed}
For models $X,Y$ put
\[
P(X,Y)\coloneqq (X\triangle Y)\cup -(X\triangle Y),
\]
the symmetric set of judgments whose problem classes receive different responses from $X$ and $Y$. For a symmetric measurable $A$, the models $X$ and $X'$ \emph{agree on $A$} if $(X\triangle X')\cap A$ is null. For models $X,Y$,
\[
d(X,Y)\coloneqq\int_{X\triangle Y}\bigl(1+\|g\|\bigr)\,d\mu(g).
\]
For a measurable $A\subseteq E$ put $\nu(A)\coloneqq2\int_A\|g\|\,d\mu(g)$.
\end{definition}

\begin{definition}[Coherence]\label{def:coherence}
A set $D\subseteq E$ is \emph{coherent} if:
\begin{enumerate}[label=(C\arabic*)]
\item $g\in D$ for every $g\in E_+\setminus\{0\}$;
\item $g\notin D$ for every $g\in E$ with $g\le0$;
\item $\lambda g\in D$ for every $g\in D$ and $\lambda>0$;
\item $g+h\in D$ for every $g,h\in D$.
\end{enumerate}
A model $X$ is \emph{almost coherent} if there is a coherent $D$ with $X\triangle D$ null. $\C$ is the set of almost coherent models.
\end{definition}

\section{The frequency measure}\label{sec:measure}

\begin{definition}[Hypotheses on $\mu$]\label{def:hypotheses}
\begin{enumerate}[label=(M\arabic*)]
\item $\mu$ is a finite symmetric Borel measure on $E$ without atoms, and $\mu(B(c,r))>0$ for every $c\in E$ and $r>0$.
\item $\int_E\|g\|\,d\mu(g)<\infty$.
\item For every measurable $A$ with $\mu(A)>0$: $\inf_{\omega\in\Omega}\int_A|g(\omega)|\,d\mu(g)>0$.
\item $\mu(\bd R)=0$ for every closed convex $R\subseteq E$ with $E_+\subseteq R$, $R+E_+\subseteq R$ and $R\cap E_{--}=\varnothing$.
\end{enumerate}
\end{definition}

\begin{lemma}\label{lem:hyperplane}
Under (M3), $\{g\in E:g(\omega)=0\}$ is null for every $\omega\in\Omega$.
\end{lemma}
\begin{proof}
Let $A\coloneqq\{g:g(\omega)=0\}$. Then $\int_A|g(\omega)|\,d\mu=0$, so $\inf_{\omega'}\int_A|g(\omega')|\,d\mu=0$. By (M3), $\mu(A)=0$.
\end{proof}

\begin{lemma}\label{lem:detecting}
There is a countable set $W\subseteq\Omega$ with $\sup_{w\in W}g(w)=\sup_{\omega\in\Omega}g(\omega)$ for every $g\in E$. Consequently, for every $g\in E$: if $g(\omega)>0$ for some $\omega\in\Omega$, then $g(w)>0$ for some $w\in W$.
\end{lemma}
\begin{proof}
Let $(v_j)_{j\in\N}$ be dense in $E$. For $j,k\in\N$ choose $\omega_{j,k}\in\Omega$ with $v_j(\omega_{j,k})>\sup_\Omega v_j-1/k$, and let $W\coloneqq\{\omega_{j,k}:j,k\in\N\}$. Let $g\in E$ and $\varepsilon>0$. Choose $j$ with $\|g-v_j\|<\varepsilon$ and $k$ with $1/k<\varepsilon$. Then
\[
\sup_Wg\ \ge\ g(\omega_{j,k})\ >\ v_j(\omega_{j,k})-\varepsilon\ >\ \sup_\Omega v_j-2\varepsilon\ \ge\ \sup_\Omega g-3\varepsilon.
\]
Since $\varepsilon$ is arbitrary and $W\subseteq\Omega$, $\sup_Wg=\sup_\Omega g$. If $g(\omega)>0$ then $\sup_Wg>0$, so $g(w)>0$ for some $w\in W$.
\end{proof}

\begin{lemma}\label{lem:extraction}
Assume (M1). Let $A_1,\dots,A_k$ be measurable and $m_1,\dots,m_k>0$ with $\sum_{i=1}^km_i\le\min_i\mu(A_i)$. Then there are pairwise disjoint measurable $B_i\subseteq A_i$ with $\mu(B_i)=m_i$.
\end{lemma}
\begin{proof}
Suppose $B_1,\dots,B_{i-1}$ have been chosen. Then $\mu\bigl(A_i\setminus(B_1\cup\dots\cup B_{i-1})\bigr)\ge\mu(A_i)-\sum_{l<i}m_l\ge m_i$. Since $\mu$ has no atoms, this set has a measurable subset $B_i$ with $\mu(B_i)=m_i$.
\end{proof}

\begin{lemma}\label{lem:separable}
Assume (M1) and (M2). Then $d$ is a finite metric on $\M$, and $(\M,d)$ is separable. Moreover, for all models $X,Y$:
\[
\mu\bigl(P(X,Y)\bigr)\le2\mu(X\triangle Y)\le2d(X,Y),\qquad \nu\bigl(P(X,Y)\bigr)\le4\,d(X,Y).
\]
\end{lemma}
\begin{proof}
Let $\lambda$ be the finite Borel measure $\lambda(A)\coloneqq\int_A(1+\|g\|)\,d\mu$; finiteness is (M2). Then $d(X,Y)=\lambda(X\triangle Y)$, and $\lambda$ and $\mu$ have the same null sets, so $d$ is a metric on $\M$. Let $\mathcal B$ be a countable base of the separable metric space $E$, and $\mathcal A$ the algebra generated by $\mathcal B$, which is countable. The class of Borel sets $A$ such that for every $\varepsilon>0$ some $A'\in\mathcal A$ has $\lambda(A\triangle A')<\varepsilon$ is a $\sigma$-algebra containing $\mathcal B$, hence contains every Borel set. Every measurable set differs from a Borel set by a null set. So $\mathcal A$ is dense in $(\M,d)$. By symmetry of $\mu$,
\begin{gather*}
\mu(P(X,Y))\le\mu(X\triangle Y)+\mu(-(X\triangle Y))=2\mu(X\triangle Y),\\
\nu(P(X,Y))\le2\int_{X\triangle Y}\|g\|\,d\mu+2\int_{-(X\triangle Y)}\|g\|\,d\mu=4\int_{X\triangle Y}\|g\|\,d\mu,
\end{gather*}
and both right-hand sides are at most $2d(X,Y)$ and $4d(X,Y)$ respectively.
\end{proof}

\begin{proposition}\label{prop:existence-measure}
There is a measure satisfying (M1)--(M4).
\end{proposition}
\begin{proof}
Let $(v_j)_{j\in\N}$ be dense in $E$, and $a_j\coloneqq c_0\,2^{-j}/(1+\|v_j\|)$ with $c_0>0$ such that $\sum_ja_j=1$. Let $N,\sigma,\zeta$ be independent random variables with $\Pr(N=j)=a_j$, $\Pr(\sigma=1)=\Pr(\sigma=-1)=\tfrac12$, and $\zeta$ real with a density $\phi$ that is continuous, everywhere positive, symmetric ($\phi(z)=\phi(-z)$), and satisfies $\int|z|\phi(z)\,dz<\infty$. Let $\mu$ be the law of $\xi\coloneqq\sigma v_N+\zeta\one$. For each $(j,s)$ the map $z\mapsto sv_j+z\one$ is continuous, so $\xi$ is a random element of $E$ and $\mu$ is a Borel probability measure.

(M1) \emph{Symmetry}: $-\xi=(-\sigma)v_N+(-\zeta)\one$, and $(N,-\sigma,-\zeta)$ has the law of $(N,\sigma,\zeta)$. \emph{No atoms}: for fixed $(j,s)$ and $h\in E$, $sv_j+z\one=h$ for at most one $z$, so $\Pr(\xi=h)=0$. A finite Borel measure on a separable metric space without point masses has no atoms: if $A$ were an atom, then for each $n$ one of countably many balls of radius $1/n$ covering $E$ would meet $A$ in a set of measure $\mu(A)$, and the intersection of these sets would be a single point of measure $\mu(A)$. \emph{Positivity on balls}: given $c\in E$ and $r>0$, choose $j$ with $\|v_j-c\|<r/2$; then $\Pr(\xi\in B(c,r))\ge\Pr(N=j,\sigma=1,|\zeta|<r/2)>0$.

(M2) $\int\|g\|\,d\mu\le\sum_ja_j\|v_j\|+\int|z|\phi(z)\,dz<\infty$.

(M3) Let $\mu(A)>0$. Since $\mu=\sum_{j,s}\tfrac{a_j}2\,\mu_{j,s}$, where $\mu_{j,s}$ is the law of $sv_j+\zeta\one$, there are $j,s$ with $\mu_{j,s}(A)>0$. Put $v\coloneqq sv_j$, $b\coloneqq a_j/2$ and $A_0\coloneqq\{z\in\R:v+z\one\in A\}$, so $\int_{A_0}\phi>0$. Define $F(\alpha)\coloneqq\int_{A_0}|\alpha+z|\phi(z)\,dz$ for $\alpha\in\R$. $F$ is finite and continuous by dominated convergence, since $|\alpha+z|\le|\alpha|+|z|$. $F(\alpha)>0$ for every $\alpha$, since the integrand is positive except at $z=-\alpha$ and $\int_{A_0}\phi>0$. Hence $\kappa\coloneqq\min_{|\alpha|\le\|v\|}F(\alpha)>0$. For every $\omega$,
\[
\int_A|g(\omega)|\,d\mu(g)\ \ge\ b\int_{A_0}|v(\omega)+z|\phi(z)\,dz\ =\ bF(v(\omega))\ \ge\ b\kappa.
\]

(M4) Let $R$ be as in (M4). Fix $v\in E$ and put $T_v\coloneqq\{z\in\R:v+z\one\in R\}$. If $z>\|v\|$ then $v+z\one\in E_{++}\subseteq R$; if $z<-\|v\|$ then $v+z\one\in E_{--}$, so $v+z\one\notin R$. $T_v$ is closed since $R$ is closed, and $z\in T_v$, $z'>z$ imply $z'\in T_v$ since $R+E_+\subseteq R$. Hence $T_v=[\alpha(v),\infty)$ for some $\alpha(v)\in\R$. If $z>\alpha(v)$ and $\|w\|<z-\alpha(v)$, then $v+z\one+w\ge v+\alpha(v)\one$, so $v+z\one+w\in R$; thus $v+z\one\in\interior R$. If $z<\alpha(v)$ then $v+z\one$ lies in the open set $E\setminus R$. So $\{z:v+z\one\in\bd R\}\subseteq\{\alpha(v)\}$, and $\mu_{j,s}(\bd R)=0$ for every $(j,s)$, because $\zeta$ has a density. Hence $\mu(\bd R)=0$.
\end{proof}

\section{Rankings and the dominance argument}\label{sec:ldg}

The definitions of this section are those of \cite{LDG}, specialised to the set of models $\M$ and to rankings that evaluate every model.

\begin{definition}[Rankings and dominance]\label{def:rankings}
A \emph{ranking} $L$ is a family $(\preceq_\omega)_{\omega\in\Omega}$ of binary relations on $\M$. $Y\preceq_\omega X$ reads: $Y$ is at least as good as $X$ at $\omega$. Its asymmetric part is
\[
Y\prec_\omega X\quad:\Longleftrightarrow\quad Y\preceq_\omega X\ \text{ and not }\ X\preceq_\omega Y .
\]
$Y$ \emph{dominates} $X$ under $L$ if $Y\preceq_\omega X$ for every $\omega\in\Omega$ and $X\not\preceq_\omega Y$ for some $\omega\in\Omega$. $X$ is \emph{undominated} under $L$ if no model dominates it. $\U_L$ is the set of undominated models and $\C_L\coloneqq\C\cap\U_L$. No property of the relations $\preceq_\omega$ is assumed.
\end{definition}

\begin{definition}[Countable selection; {\cite[Definitions~3.4 and~4.3]{LDG}}]\label{def:selection}
$\Delta$ is the set of finitely additive probabilities on $\Omega$ with the topology of \cite{LDG}, under which it is compact; $p_\omega\coloneqq p(\{\omega\})$, and $p$ is \emph{finitely supported} if $p(F)=1$ for a finite $F\subseteq\Omega$. A ranking $L$ satisfies \emph{countable selection on $\mathcal I\subseteq\M$} if there are models $c^j_p\in\M$, $j\in\N$, $p\in\Delta$, such that for every $x\in\mathcal I$ there are $j\in\N$ and closed sets $A_\omega(x)\subseteq\Delta$, $\omega\in\Omega$, with
\begin{align}
&\forall\omega\in\Omega\ \forall p\in A_\omega(x):\ c^j_p\preceq_\omega x;\tag{F1}\\
&\forall p\in\Delta\text{ finitely supported}\ \exists\omega\in\Omega:\ p_\omega>0\text{ and }p\in A_\omega(x).\tag{F2}
\end{align}
For a family $\Lf$ of rankings:
\begin{itemize}[leftmargin=3em]
\item[(E)] every model is evaluated by some $L\in\Lf$ (here every ranking evaluates every model);
\item[(A)] $\C=\bigcup_{L\in\Lf}\C_L$;
\item[(S)] every $L\in\Lf$ satisfies countable selection on $\M\setminus\C$;
\item[(T)] for every $L\in\Lf$ the models $c^j_p$ in (S) can be chosen so that each $x\in\M\setminus\C$ has $j$ and closed sets $A_\omega(x)$ satisfying (F1), (F2) and
\begin{equation}
\forall p\in\textstyle\bigcap_{\omega\in\Omega}A_\omega(x)\ \exists\omega\in\Omega:\ x\not\preceq_\omega c^j_p.\tag{F3}
\end{equation}
\end{itemize}
A model $m$ is \emph{rationally permissible} relative to $\Lf$ if $m\in\U_L$ for some $L\in\Lf$, and \emph{irrational} otherwise.
\end{definition}

In \cite{LDG} the class $\C$ is an arbitrary class of models fixed in advance. Definition~\ref{def:selection} and the following theorem are applied below to $\C$ and, in Section~\ref{sec:convex}, to a class $\C^0$ in place of $\C$, with $\C^0_L\coloneqq\C^0\cap\U_L$.

\begin{theorem}[{\cite[Theorem~4.4(d)]{LDG}}]\label{thm:ldg}
Let $\C^\ast\subseteq\M$ be any class of models, and read (A), (S) and (T) with $\C^\ast$ in place of $\C$. If a family $\Lf$ of rankings satisfies (A), (S) and (T), then the rationally permissible models are exactly the members of $\C^\ast$, and the irrational models are exactly the members of $\M\setminus\C^\ast$.
\end{theorem}

\section{Axioms}\label{sec:axioms}

Throughout, $t>1$ is fixed. All set equalities and inclusions between models and between sets of judgments are understood up to null sets.

\begin{definition}[Required comparisons of responses]\label{def:required}
The \emph{required weak transitions} are the ordered pairs $(b,a)$ of distinct responses with $a\rhd b$, read: a change from response $b$ to response $a$. The \emph{required strict transitions} are the required weak transitions other than $(\con,\wro)$.
\end{definition}

\begin{axiom}{G (Guidance)}
For all models $X,Y$ and every $\omega\in\Omega$:
\begin{enumerate}[label=(\roman*)]
\item if, for almost every $g\in P(X,Y)$ with $g(\omega)\ne0$, $(\rho_\omega(X,g),\rho_\omega(Y,g))$ is a required weak transition, then $Y\preceq_\omega X$;
\item if moreover $\mu(P(X,Y))>0$ and, for almost every $g\in P(X,Y)$ with $g(\omega)\ne0$, $(\rho_\omega(X,g),\rho_\omega(Y,g))$ is a required strict transition, then $Y\prec_\omega X$.
\end{enumerate}
\end{axiom}
\noindent At a given world, a change that moves responses up the order correct, silent, wrong, confused is at least as good there, and strictly better if no changed response moves only from confused to wrong.

\begin{definition}[Averaging instance]\label{def:AV}
$(X,Z)$ is an \emph{averaging instance} if there are $f_1,f_2,h\in E$, $r>0$, $\varepsilon>0$, $m>0$ and measurable sets $G_1,G_2,H$ with
\begin{gather*}
G_1\subseteq U(X)\cap B(f_1,r),\quad G_2\subseteq U(X)\cap B(f_2,r),\quad G_1\cap G_2=\varnothing,\\
H\subseteq S(X)\cap B(h,r),\quad \mu(G_1)=\mu(G_2)=m,\quad \mu(H)=2m,\\
r<\|h\|,\quad h\ge\tfrac12(f_1+f_2)+(\varepsilon+2r)\one,
\end{gather*}
and $Z=(X\setminus(G_1\cup G_2))\cup H$.
\end{definition}

\begin{axiom}{AV (Averaging)}
If $(X,Z)$ is an averaging instance, then $Z\prec_\omega X$ for every $\omega\in\Omega$.
\end{axiom}
\noindent Replacing two equally frequent sets of recommendations near $f_1$ and $f_2$ by recommendations of twice that frequency, on silent problems, near a gamble uniformly above their average makes the advice strictly better at every world.

\begin{definition}[Weighted sure-loss withdrawal instance]\label{def:WSL}
$(X,Z)$ is a \emph{weighted sure-loss withdrawal instance} if there are $k\in\N$, $f_1,\dots,f_k\in E$, $r>0$, $\varepsilon>0$, $m_1,\dots,m_k>0$ and pairwise disjoint measurable sets $G_1,\dots,G_k$ with
\[
G_i\subseteq U(X)\cap B(f_i,r),\quad \mu(G_i)=m_i,\quad \sum_{i=1}^km_if_i\le-\Bigl(\varepsilon+r\sum_{i=1}^km_i\Bigr)\one,
\]
and $Z=X\setminus(G_1\cup\dots\cup G_k)$.
\end{definition}

\begin{axiom}{WSL (Weighted sure-loss withdrawal)}
If $(X,Z)$ is a weighted sure-loss withdrawal instance, then $Z\prec_\omega X$ for every $\omega\in\Omega$.
\end{axiom}
\noindent Withdrawing sets of recommendations whose frequency-weighted sum is uniformly negative makes the advice strictly better at every world.

\begin{definition}[Concentration instance]\label{def:CONC}
$(X,Z)$ is a \emph{concentration instance} (for the fixed $t>1$) if there are $f,h\in E$, $r>0$, $\varepsilon>0$, $m>0$ and measurable sets $G,H$ with
\begin{gather*}
G\subseteq U(X)\cap B(f,r),\quad H\subseteq S(X)\cap B(h,r),\quad \mu(G)=m,\quad \mu(H)=m/t,\\
r<\|h\|,\quad h\ge tf+\bigl(\varepsilon+(t+1)r\bigr)\one,
\end{gather*}
and $Z=(X\setminus G)\cup H$.
\end{definition}

\begin{axiom}{CONC (Concentration)}
If $(X,Z)$ is a concentration instance, then $Z\prec_\omega X$ for every $\omega\in\Omega$.
\end{axiom}
\noindent Replacing recommendations near $f$ by recommendations of $1/t$ times the frequency, on silent problems, near a gamble uniformly above $tf$ makes the advice strictly better at every world.

\begin{axiom}{I (Independence)}
For all models $X,Y,X',Y'$: if $P(X,Y)=P(X',Y')\eqqcolon P$, $X$ and $X'$ agree on $P$, and $Y$ and $Y'$ agree on $P$, then for every $\omega\in\Omega$
\[
Y\preceq_\omega X\iff Y'\preceq_\omega X'\qquad\text{and}\qquad X\preceq_\omega Y\iff X'\preceq_\omega Y'.
\]
\end{axiom}
\noindent The comparison of two models at a world depends only on which problem classes they treat differently and on how each treats those classes.

\begin{definition}[Required improvements]\label{def:J}
For models $X,Z$:
\begin{enumerate}[label=(\alph*)]
\item $J_\GU(X,Z)$ holds if $\mu(P(X,Z))>0$ and, for every $\omega\in\Omega$ and almost every $g\in P(X,Z)$ with $g(\omega)\ne0$, $(\rho_\omega(X,g),\rho_\omega(Z,g))$ is a required strict transition;
\item $J(X,Z)$ holds if $J_\GU(X,Z)$ holds or $(X,Z)$ is an averaging, a weighted sure-loss withdrawal or a concentration instance.
\end{enumerate}
\end{definition}

\begin{axiom}{UR (Uniform robustness)}
For all models $X,Z$ with $J(X,Z)$ there is $\delta>0$ such that for every model $Y$:
\[
d(Y,Z)<\delta\ \Longrightarrow\ Y\prec_\omega X\ \text{ for every }\omega\in\Omega.
\]
\end{axiom}
\noindent If a change is required by Guidance at every world, or by Averaging, Weighted sure-loss withdrawal or Concentration, then every sufficiently small alteration of its result is still strictly better at every world, with one bound on the size of the alteration for all worlds.

\begin{definition}\label{def:class}
$\K_t$ is the class of rankings satisfying Axioms G, AV, WSL, CONC, I and UR.
\end{definition}

\section{Classification}\label{sec:classification}

\begin{theorem}\label{thm:classification}
Assume (M1)--(M4). If $X\in\M$ and there is no $Z\in\M$ with $J(X,Z)$, then $X\in\C$.
\end{theorem}

\begin{proof}
Write $U\coloneqq U(X)$, $S\coloneqq S(X)$, $K\coloneqq K(X)$. For a measurable $A$ let $\mu_A(B)\coloneqq\mu(A\cap B)$, and let $\supp\mu_A$ be the set of $c\in E$ with $\mu(A\cap B(c,r))>0$ for every $r>0$. Since $E$ is separable, $\supp\mu_A$ is closed and $\mu(A\setminus\supp\mu_A)=0$. Put $R\coloneqq\supp\mu_U$. By Lemma~\ref{lem:hyperplane}, for each $\omega$ the set $\{g:g(\omega)=0\}$ is null; this is used without mention when responses at $\omega$ are compared.

\emph{Step 0.}
\begin{enumerate}[label=(\alph*)]
\item $K$ is null. Otherwise let $Z\coloneqq X\setminus K$. Then $P(X,Z)=K$ is not null, and for every $\omega$ and $g\in K$ with $g(\omega)\ne0$, $\rho_\omega(X,g)=\con$ and $\rho_\omega(Z,g)=\sil$. So $J_\GU(X,Z)$, a contradiction. Hence $X=U$.
\item $E_{++}\setminus U$ is null. Otherwise put $A\coloneqq E_{++}\setminus U$ and $Z\coloneqq(X\setminus(-A))\cup A$. $A\cap(-A)=\varnothing$ since $-A\subseteq E_{--}$, and $P(X,Z)\supseteq A$ has positive measure. For $g\in A$: $g\in Z$ and $-g\notin Z$, and $g^\omega=g$ for every $\omega$, so $\rho_\omega(Z,g)=\cor$, while $\rho_\omega(X,g)\in\{\sil,\wro\}$ for almost every $g\in A$ by (a). So $J_\GU(X,Z)$, a contradiction.
\item $U\cap E_{--}$ is null. Otherwise put $A\coloneqq U\cap E_{--}$ and $Z\coloneqq X\setminus A$. For $g\in A$ and every $\omega$, $\rho_\omega(X,g)=\wro$ and $\rho_\omega(Z,g)=\sil$. So $J_\GU(X,Z)$, a contradiction.
\item $E_+\subseteq R$, in particular $0\in R$, and $R\cap E_{--}=\varnothing$. For $c\in E_{++}$ and $r>0$, $B(c,r)\cap E_{++}$ is a nonempty open set, so it has positive measure by (M1), and by (b) $\mu(U\cap B(c,r))>0$; thus $E_{++}\subseteq R$. For $c\in E_+$, $c+\tfrac1n\one\in E_{++}$ and $c+\tfrac1n\one\to c$; $R$ is closed, so $c\in R$. $E_{--}$ is open and $\mu(U\cap E_{--})=0$ by (c), so no point of $E_{--}$ is in $R$.
\end{enumerate}

\emph{Step 1. If $f_1,f_2\in R$, $\eta>0$ and $h\ge\frac12(f_1+f_2)+\eta\one$, then $h\in R$.} If $h=0$, then $h\in R$ by Step 0(d). Let $h\ne0$ and suppose $h\notin R$. Choose $r_0>0$ with $\mu(U\cap B(h,r_0))=0$, and $r$ with $0<r<\min(r_0,\|h\|)$ and $r\le\eta/4$. Put $\varepsilon\coloneqq\eta/2$; then $h\ge\frac12(f_1+f_2)+(\varepsilon+2r)\one$. Let $V\coloneqq B(h,r)$. By (M1), $\mu(V)>0$, and by Step 0(a) almost every $g\in V$ lies in $S$ or satisfies $-g\in U$. Since $f_1,f_2\in R$, $\mu(U\cap B(f_i,r))>0$.
\begin{itemize}
\item If $\mu(S\cap V)>0$: choose $m>0$ with $2m\le\min\{\mu(U\cap B(f_1,r)),\mu(U\cap B(f_2,r)),\mu(S\cap V)\}$. By Lemma~\ref{lem:extraction} there are disjoint $G_1\subseteq U\cap B(f_1,r)$, $G_2\subseteq U\cap B(f_2,r)$ with $\mu(G_i)=m$, and $H\subseteq S\cap V$ with $\mu(H)=2m$; $H$ is disjoint from $G_1\cup G_2\subseteq U$. Then $(X,(X\setminus(G_1\cup G_2))\cup H)$ is an averaging instance, a contradiction.
\item If $\mu(S\cap V)=0$: then $-g\in U$ for almost every $g\in V$, so by symmetry $\mu(U\cap B(-h,r))=\mu(V)>0$. Choose $m>0$ with $4m\le\min\{\mu(U\cap B(f_1,r)),\mu(U\cap B(f_2,r)),\mu(U\cap B(-h,r))\}$ and, by Lemma~\ref{lem:extraction}, pairwise disjoint $G_1\subseteq U\cap B(f_1,r)$, $G_2\subseteq U\cap B(f_2,r)$, $G_3,G_4\subseteq U\cap B(-h,r)$ of measure $m$ each. With centres $f_1,f_2,-h,-h$:
\[
m f_1+mf_2-2mh\le-2m\eta\one\le-\bigl(m\eta+4mr\bigr)\one,
\]
since $4r\le\eta$. So $(X,X\setminus(G_1\cup\dots\cup G_4))$ is a weighted sure-loss withdrawal instance with $\varepsilon\coloneqq m\eta$, a contradiction.
\end{itemize}
Hence $h\in R$.

\emph{Consequences of Step 1.} (i) If $f_1,f_2\in R$ then $\frac12(f_1+f_2)+\frac1n\one\in R$ for all $n$, so $\frac12(f_1+f_2)\in R$ because $R$ is closed. A closed set closed under midpoints is convex: the dyadic convex combinations of two points of $R$ lie in $R$ by induction, and they are dense in the segment. (ii) If $f\in R$ and $e\in E_+$, then Step 1 with $f_1=f_2=f$ and $h=f+e+\frac1n\one$ gives $f+e+\frac1n\one\in R$, so $f+e\in R$. Thus $R+E_+\subseteq R$.

\emph{Step 2. If $f\in R$ then $tf\in R$.} Let $\eta>0$ and $h\coloneqq tf+\eta\one$. If $h=0$ then $f=-(\eta/t)\one\in E_{--}$, contradicting Step 0(d); so $h\ne0$. Suppose $h\notin R$. Choose $r_0$ with $\mu(U\cap B(h,r_0))=0$, and $r$ with $0<r<\min(r_0,\|h\|)$ and $(t+1)r\le\eta/2$. Put $\varepsilon\coloneqq\eta/2$; then $h\ge tf+(\varepsilon+(t+1)r)\one$. Let $V\coloneqq B(h,r)$.
\begin{itemize}
\item If $\mu(S\cap V)>0$: choose $m>0$ with $m\le\mu(U\cap B(f,r))$ and $m/t\le\mu(S\cap V)$, and $G\subseteq U\cap B(f,r)$, $H\subseteq S\cap V$ with $\mu(G)=m$, $\mu(H)=m/t$. Then $(X,(X\setminus G)\cup H)$ is a concentration instance, a contradiction.
\item If $\mu(S\cap V)=0$: as in Step 1, $\mu(U\cap B(-h,r))>0$. Choose $m>0$ with $m+m/t\le\min\{\mu(U\cap B(f,r)),\mu(U\cap B(-h,r))\}$ and, by Lemma~\ref{lem:extraction}, disjoint $G_1\subseteq U\cap B(f,r)$, $G_2\subseteq U\cap B(-h,r)$ with $\mu(G_1)=m$, $\mu(G_2)=m/t$. With centres $f,-h$:
\[
mf-\tfrac mt h=\tfrac mt(tf-h)\le-\tfrac mt\eta\one\le-\Bigl(\tfrac{m\eta}{2t}+\bigl(m+\tfrac mt\bigr)r\Bigr)\one,
\]
since $(m+m/t)r=\frac mt(t+1)r\le\frac{m\eta}{2t}$. So $(X,X\setminus(G_1\cup G_2))$ is a weighted sure-loss withdrawal instance with $\varepsilon\coloneqq m\eta/(2t)$, a contradiction.
\end{itemize}
Hence $tf+\eta\one\in R$ for every $\eta>0$, and $tf\in R$ since $R$ is closed.

\emph{Step 3. $R$ is a closed convex cone with $E_+\subseteq R$, $R+E_+\subseteq R$ and $R\cap E_{--}=\varnothing$.} By Step 2, $t^jR\subseteq R$ for all $j\in\N$. Let $f\in R$ and $\lambda>0$; choose $j$ with $t^j\ge\lambda$. Then $\lambda f=\frac{\lambda}{t^j}(t^jf)+\bigl(1-\frac{\lambda}{t^j}\bigr)0\in R$ by convexity and $0\in R$. The remaining properties are Step 0(d) and Consequence (ii).

\emph{Step 4. $U\setminus\interior R$ is null.} $\interior R\supseteq E_{++}\ne\varnothing$. By Step 3 and (M4), $\mu(\bd R)=0$. Hence $\mu(U\setminus\interior R)\le\mu(U\setminus R)+\mu(\bd R)=0$.

\emph{Step 5. $\interior R\setminus U$ is null.} Suppose $Q\coloneqq\interior R\setminus U$ has positive measure. Since $\mu(Q\setminus\supp\mu_Q)=0$, there is $h\in Q\cap\supp\mu_Q$. Then $h\ne0$: if $0\in\interior R$ then $-\tau\one\in R$ for small $\tau>0$, contradicting $R\cap E_{--}=\varnothing$. Choose $\rho$ with $0<\rho<\|h\|$ and $B(h,\rho)\subseteq\interior R$. Put $r\coloneqq\varepsilon\coloneqq\rho/4$ and $f\coloneqq h-(\varepsilon+2r)\one$. Then $\|f-h\|=3\rho/4$, so $f\in R$ and $\mu(U\cap B(f,r))>0$. Let $V\coloneqq B(h,r)$; then $\mu(V\setminus U)\ge\mu(Q\cap V)>0$, and $h=\frac12(f+f)+(\varepsilon+2r)\one$.
\begin{itemize}
\item If $\mu(S\cap V)>0$: by Lemma~\ref{lem:extraction} choose disjoint $G_1,G_2\subseteq U\cap B(f,r)$ of measure $m$ and $H\subseteq S\cap V$ of measure $2m$, for small $m>0$. This is an averaging instance with $f_1=f_2=f$, a contradiction.
\item If $\mu(S\cap V)=0$: by Step 0(a), $-g\in U$ for almost every $g\in V\setminus U$, so $\mu(U\cap B(-h,r))>0$. Choose pairwise disjoint $G_1,G_2\subseteq U\cap B(f,r)$ and $G_3,G_4\subseteq U\cap B(-h,r)$ of measure $m$. Then $2mf-2mh=-2m(\varepsilon+2r)\one\le-(m\varepsilon+4mr)\one$, a weighted sure-loss withdrawal instance, a contradiction.
\end{itemize}

\emph{Step 6.} Let $D\coloneqq\interior R\cup(E_+\setminus\{0\})$.
\begin{itemize}
\item (C1) holds by definition.
\item (C2): if $g\le0$ and $g\in\interior R$, then $g-\varepsilon\one\in R$ for small $\varepsilon>0$, and $g-\varepsilon\one\in E_{--}$, contradicting Step 3. No $g\in E_+\setminus\{0\}$ satisfies $g\le0$.
\item (C3): for $\lambda>0$, $g\mapsto\lambda g$ is a homeomorphism of $E$ mapping $R$ onto $R$ (Step 3), hence $\interior R$ onto $\interior R$; and $\lambda(E_+\setminus\{0\})=E_+\setminus\{0\}$.
\item (C4): if $g,h\in\interior R$ then $\frac12(g+h)\in\interior R$, since $\interior R$ is convex, and $g+h=2\cdot\frac12(g+h)\in\interior R$ by (C3). If $g\in\interior R$, $e\in E_+$ and $B(g,\varepsilon)\subseteq R$, then $B(g+e,\varepsilon)=B(g,\varepsilon)+e\subseteq R+E_+\subseteq R$, so $g+e\in\interior R$. $E_+\setminus\{0\}$ is closed under addition.
\end{itemize}
So $D$ is coherent. Finally, $X=U$ (Step 0(a)), $U\triangle\interior R$ is null (Steps 4, 5), and $D\setminus\interior R\subseteq E_+\setminus\interior R\subseteq R\setminus\interior R=\bd R$ is null. Hence $X\triangle D$ is null and $X\in\C$.
\end{proof}

\begin{corollary}\label{cor:direct}
Assume (M1)--(M4). Let $L\in\K_t$ and $X\in\M\setminus\C$. Then there is $Z\in\M$ with $J(X,Z)$, and $Z\prec_\omega X$ for every $\omega\in\Omega$; in particular $Z$ dominates $X$ under $L$.
\end{corollary}
\begin{proof}
By Theorem~\ref{thm:classification} there is $Z$ with $J(X,Z)$. If $J_\GU(X,Z)$, Axiom G(ii) gives $Z\prec_\omega X$ for every $\omega$. Otherwise Axiom AV, WSL or CONC does. Since $\Omega\ne\varnothing$, $Z$ dominates $X$.
\end{proof}

\begin{lemma}\label{lem:noJ}
Assume (M1)--(M3). If $T\in\C$, there is no $Z\in\M$ with $J(T,Z)$.
\end{lemma}

\begin{proof}
Let $D$ be coherent with $T\triangle D$ null. By (C2) and (C4), $g\in D$ and $-g\in D$ cannot both hold, since $g+(-g)=0\le0$. Up to null sets, $U(T)=D$ and $K(T)=\varnothing$.

\emph{Averaging.} Suppose $(T,Z)$ is an averaging instance with data as in Definition~\ref{def:AV}. Since $\mu(G_i)>0$ and $G_i\subseteq U(T)$, there are $f_1'\in G_1\cap D$ and $f_2'\in G_2\cap D$. For $h'\in H$: $h'\ge h-r\one\ge\frac12(f_1+f_2)+(\varepsilon+r)\one\ge\frac12(f_1'+f_2')+\varepsilon\one$, because $f_i'\le f_i+r\one$. By (C3), (C4) and (C1), $h'=\frac12f_1'+\frac12f_2'+(h'-\frac12(f_1'+f_2'))\in D$. So $H\subseteq D$, and $H\subseteq S(T)$ forces $\mu(H)=0$, a contradiction.

\emph{Weighted sure-loss withdrawal.} With data as in Definition~\ref{def:WSL}, choose $f_i'\in G_i\cap D$. Then $\sum_im_if_i'\le\sum_im_i(f_i+r\one)\le-\varepsilon\one$, so $\sum_im_if_i'\le0$, while $\sum_im_if_i'\in D$ by (C3), (C4). This contradicts (C2).

\emph{Concentration.} With data as in Definition~\ref{def:CONC}, choose $f'\in G\cap D$. For $h'\in H$: $h'\ge h-r\one\ge tf+(\varepsilon+tr)\one\ge tf'+\varepsilon\one$, so $h'=tf'+(h'-tf')\in D$. Thus $H\subseteq D$, a contradiction as before.

\emph{$J_\GU$.} Suppose $J_\GU(T,Z)$, and let $W$ be as in Lemma~\ref{lem:detecting}. For each $w\in W$ let $N_w$ be the null set of $g\in P(T,Z)$ for which $g(w)=0$ or $(\rho_w(T,g),\rho_w(Z,g))$ is not a required strict transition. Let $N\coloneqq\bigcup_{w\in W}N_w\cup\bigl((T\triangle D)\cup-(T\triangle D)\bigr)$, a null set. Since $\mu(P(T,Z))>0$, there is $g\in P(T,Z)\setminus N$, with $g\ne0$. Then $T\cap[g]=D\cap[g]$, and the response of $T$ to $[g]$ strictly improves at every $w\in W$.
\begin{itemize}
\item If $D\cap[g]=\{g'\}$: by (C2), $g'(\omega)>0$ for some $\omega$, so $g'(w)>0$ for some $w\in W$ (Lemma~\ref{lem:detecting}). Then $\rho_w(T,g)=\cor$, which is maximal for $\rhd$; no change of it is a required transition.
\item If $D\cap[g]=\varnothing$: by (C1), neither $g$ nor $-g$ lies in $E_+\setminus\{0\}$, and $g\ne0$, so $g$ takes a strictly positive and a strictly negative value. By Lemma~\ref{lem:detecting} there are $w_1,w_2\in W$ with $g(w_1)>0>g(w_2)$. $\rho_w(T,g)=\sil$ for every $w$. If $Z\cap[g]=\{g\}$ then $\rho_{w_2}(Z,g)=\wro$; if $Z\cap[g]=\{-g\}$ then $\rho_{w_1}(Z,g)=\wro$; if $Z\cap[g]=[g]$ then $\rho_{w_1}(Z,g)=\con$. In each case some $w\in W$ has $\sil\rhd\rho_w(Z,g)$.
\end{itemize}
Both cases contradict the choice of $g$.
\end{proof}

\section{Countable selection}\label{sec:selection}

\begin{lemma}\label{lem:general-selection}
Let $(\M',d')$ be a separable metric space, $(\preceq_\omega)_{\omega\in\Omega}$ binary relations on $\M'$, $\C'\subseteq\M'$, and $J'\subseteq\M'\times\M'$. Assume:
\begin{enumerate}[label=(\roman*)]
\item for every $x\in\M'\setminus\C'$ there is $z$ with $J'(x,z)$;
\item for all $x,z$ with $J'(x,z)$ there is $\delta>0$ such that every $y$ with $d'(y,z)<\delta$ satisfies $y\prec_\omega x$ for every $\omega$.
\end{enumerate}
Let $(d_j)_{j\in\N}$ be a sequence dense in $\M'$; it does not depend on $x$. Put $c^j_p\coloneqq d_j$ for all $j\in\N$, $p\in\Delta$. Then for every $x\in\M'\setminus\C'$ there is $j$ such that, with $A_\omega(x)\coloneqq\Delta$ for every $\omega$, (F1), (F2) and (F3) hold, and $x\not\preceq_\omega c^j_p$ for all $\omega\in\Omega$ and $p\in\Delta$.
\end{lemma}
\begin{proof}
Let $x\notin\C'$. By (i) and (ii) there are $z$ and $\delta>0$ such that $B'\coloneqq\{y:d'(y,z)<\delta\}$ satisfies $y\prec_\omega x$ for all $y\in B'$ and $\omega\in\Omega$. $B'$ is a nonempty open set, so some $d_j\in B'$; fix the least such $j$. Then $d_j\preceq_\omega x$ and $x\not\preceq_\omega d_j$ for every $\omega$. $\Delta$ is closed in $\Delta$. (F1): $c^j_p=d_j\preceq_\omega x$ for all $\omega$, $p$. (F2): if $p(F)=1$ with $F$ finite, then $\sum_{\omega\in F}p_\omega=1$ by finite additivity, so $p_\omega>0$ for some $\omega$, and $p\in\Delta=A_\omega(x)$. (F3): for every $p\in\Delta$ and every $\omega$, $x\not\preceq_\omega c^j_p$; $\Omega\ne\varnothing$.
\end{proof}

\begin{theorem}\label{thm:selection}
Assume (M1)--(M4). Every $L\in\K_t$ satisfies countable selection on $\M\setminus\C$ with models $c^j_p$ satisfying (F3), independently of the choice of $L$: the models $c^j_p\coloneqq d_j$, for any sequence $(d_j)$ dense in $(\M,d)$, serve for every $L\in\K_t$.
\end{theorem}
\begin{proof}
$(\M,d)$ is a separable metric space (Lemma~\ref{lem:separable}). Condition (i) of Lemma~\ref{lem:general-selection}, with $\C'\coloneqq\C$ and $J'\coloneqq J$, is Theorem~\ref{thm:classification}. Condition (ii) is Axiom UR. Apply Lemma~\ref{lem:general-selection}.
\end{proof}

\section{A member of the class that leaves every almost coherent model undominated}\label{sec:existence}

In this section (M1)--(M4) are assumed. The functions $\Phi_\omega$, $\Gamma_\omega$ and $\gamma$ are used only to define and analyse the ranking $L^\ast$ below.

\begin{definition}\label{def:Phi}
For $X\in\M$, $\omega\in\Omega$ and $g\in E$ put
\[
\varphi_\omega(g;X)\coloneqq g(\omega)\one_{U(X)}(g)-|g(\omega)|\one_{K(X)}(g),\qquad \Phi_\omega(X)\coloneqq\int_E\varphi_\omega(g;X)\,d\mu(g),
\]
which is finite by (M2). For $X,Z\in\M$ put $\Gamma_\omega(X,Z)\coloneqq\Phi_\omega(Z)-\Phi_\omega(X)$ and $\gamma(X,Z)\coloneqq\inf_{\omega\in\Omega}\Gamma_\omega(X,Z)$.
\end{definition}

\begin{lemma}\label{lem:Phi}
Let $X,Y\in\M$ and $\omega\in\Omega$.
\begin{enumerate}[label=(\alph*)]
\item For every symmetric measurable $A\supseteq P(X,Y)$,
\[
\Phi_\omega(Y)-\Phi_\omega(X)=\int_A\bigl(\varphi_\omega(g;Y)-\varphi_\omega(g;X)\bigr)d\mu(g)=\tfrac12\int_A\bigl(v_\omega(g;Y)-v_\omega(g;X)\bigr)d\mu(g),
\]
where $v_\omega(g;X)\coloneqq\varphi_\omega(g;X)+\varphi_\omega(-g;X)$ equals $|g(\omega)|$, $0$, $-|g(\omega)|$, $-2|g(\omega)|$ according as $\rho_\omega(X,g)$ is $\cor$, $\sil$, $\wro$, $\con$ (for $g(\omega)\ne0$).
\item $|\Phi_\omega(Y)-\Phi_\omega(X)|\le\nu(P(X,Y))$.
\item If $\mu(P(X,Y))>0$ and $\rho_\omega(Y,g)\rhd\rho_\omega(X,g)$ for almost every $g\in P(X,Y)$ with $g(\omega)\ne0$, then
\[
\Phi_\omega(Y)-\Phi_\omega(X)\ge\tfrac12\int_{P(X,Y)}|g(\omega)|\,d\mu(g)>0.
\]
\item If $J_\GU(X,Y)$, then $\gamma(X,Y)\ge\frac12\inf_{\omega}\int_{P(X,Y)}|g(\omega)|\,d\mu(g)>0$.
\item If $(X,Y)$ is an averaging, weighted sure-loss withdrawal or concentration instance with data as in Definitions~\ref{def:AV}--\ref{def:CONC}, then $\gamma(X,Y)\ge 2m\varepsilon$, $\varepsilon$, $m\varepsilon/t$ respectively. In particular $J(X,Y)$ implies $\gamma(X,Y)>0$.
\end{enumerate}
\end{lemma}

\begin{proof}
(a) If $g\notin P(X,Y)$ then $g$ and $-g$ each lie in both or neither of $X,Y$, so $\varphi_\omega(g;X)=\varphi_\omega(g;Y)$; this gives the first equality. Since $A$ and $\mu$ are symmetric, $\int_A\varphi_\omega(-g;\cdot)\,d\mu=\int_A\varphi_\omega(g;\cdot)\,d\mu$, which gives the second. For the values of $v_\omega$, let $g^\omega$ be as in Definition~\ref{def:responses} and $s\coloneqq|g(\omega)|=g^\omega(\omega)$. $v_\omega(g;X)=\varphi_\omega(g^\omega;X)+\varphi_\omega(-g^\omega;X)$. If $\rho_\omega(X,g)=\cor$: $g^\omega\in U(X)$, $-g^\omega\notin X$, so $v=s$. If $\sil$: $v=0$. If $\wro$: $-g^\omega\in U(X)$, $g^\omega\notin X$, so $v=-g^\omega(\omega)=-s$. If $\con$: both lie in $K(X)$, so $v=-2s$.

(b) For every $g$, $\varphi_\omega(g;\cdot)\in[-|g(\omega)|,|g(\omega)|]$, so by (a) with $A=P(X,Y)$, $|\Phi_\omega(Y)-\Phi_\omega(X)|\le\int_{P(X,Y)}2|g(\omega)|\,d\mu\le\nu(P(X,Y))$.

(c) By the values in (a), each step up in $\rhd$ raises $v_\omega$ by at least $|g(\omega)|$. Apply (a) with $A=P(X,Y)$. The integral is positive because $|g(\omega)|>0$ almost everywhere (Lemma~\ref{lem:hyperplane}) and $\mu(P(X,Y))>0$.

(d) Each required strict transition is a step up in $\rhd$. Apply (c) at every $\omega$, and (M3) to $A=P(X,Y)$.

(e) Consider an averaging instance. Since $G_i\subseteq U(X)$, $-g\notin X$ for $g\in G_i$. Since $H\subseteq S(X)$ and $r<\|h\|$, we have $H\cap(-H)=\varnothing$ (if $g,-g\in B(h,r)$ then $2\|h\|\le\|h-g\|+\|h+g\|<2r$), $H\cap(-G_i)=\varnothing$ (if $-g\in G_i$ then $-g\in X$, so $g\notin S(X)$) and $H\cap G_i=\varnothing$. Hence for $g\in G_1\cup G_2$: $\varphi_\omega(g;X)=g(\omega)$, $\varphi_\omega(g;Y)=0$; for $g\in H$: $\varphi_\omega(g;X)=0$, $\varphi_\omega(g;Y)=g(\omega)$; and $\varphi_\omega(g;X)=\varphi_\omega(g;Y)=0$ for $g\in-(G_1\cup G_2\cup H)$. By (a),
\begin{align*}
\Gamma_\omega(X,Y)&=\int_Hg(\omega)\,d\mu-\int_{G_1}g(\omega)\,d\mu-\int_{G_2}g(\omega)\,d\mu\\
&\ge2m(h(\omega)-r)-m(f_1(\omega)+r)-m(f_2(\omega)+r)\ge2m\varepsilon,
\end{align*}
using $|g(\omega)-c(\omega)|\le\|g-c\|$ and the inequality of Definition~\ref{def:AV}. For a weighted sure-loss withdrawal instance, similarly $\Gamma_\omega=-\sum_i\int_{G_i}g(\omega)\,d\mu\ge-\sum_im_i(f_i(\omega)+r)\ge\varepsilon$. For a concentration instance, $\Gamma_\omega=\int_Hg(\omega)\,d\mu-\int_Gg(\omega)\,d\mu\ge\frac mt(h(\omega)-r)-m(f(\omega)+r)\ge\frac mt\bigl(tf(\omega)+\varepsilon+tr\bigr)-mf(\omega)-mr=\frac{m\varepsilon}t$. These bounds hold for every $\omega$. The last sentence follows with (d).
\end{proof}

\begin{definition}[The ranking $L^\ast$]\label{def:R}
For models $X,Y$ and $\omega\in\Omega$, let $Y\preceq^\ast_\omega X$ hold if and only if one of the following holds:
\begin{enumerate}[label=(L\arabic*)]
\item $X\triangle Y$ is null;
\item $\mu(P(X,Y))>0$ and $\rho_\omega(Y,g)\rhd\rho_\omega(X,g)$ for almost every $g\in P(X,Y)$ with $g(\omega)\ne0$;
\item there is $Z\in\M$ such that $J(X,Z)$, $P(X,Z)\subseteq P(X,Y)$ and $\nu(P(Y,Z))<\frac14\gamma(X,Z)$.
\end{enumerate}
Condition (L3) does not depend on $\omega$. $L^\ast$ denotes the ranking $(\preceq^\ast_\omega)_{\omega\in\Omega}$.
\end{definition}

\begin{lemma}\label{lem:monotone}
For all $X,Y\in\M$ and $\omega\in\Omega$: if $Y\preceq^\ast_\omega X$ then $\Phi_\omega(Y)\ge\Phi_\omega(X)$; if (L2) or (L3) holds then $\Phi_\omega(Y)>\Phi_\omega(X)$. Consequently $Y\prec^\ast_\omega X$ if and only if $Y\preceq^\ast_\omega X$ and $X\triangle Y$ is not null.
\end{lemma}
\begin{proof}
Under (L1), $\Phi_\omega(Y)=\Phi_\omega(X)$. Under (L2), Lemma~\ref{lem:Phi}(c). Under (L3), by Lemma~\ref{lem:Phi}(b) and~(e),
\begin{align*}
\Phi_\omega(Y)-\Phi_\omega(X)&=\Gamma_\omega(X,Z)+\bigl(\Phi_\omega(Y)-\Phi_\omega(Z)\bigr)\\
&\ge\gamma(X,Z)-\nu(P(Y,Z))>\tfrac34\gamma(X,Z)>0.
\end{align*}
If $Y\preceq^\ast_\omega X$ and $X\triangle Y$ is not null, then (L2) or (L3) holds, so $\Phi_\omega(Y)>\Phi_\omega(X)$; then $X\preceq^\ast_\omega Y$ would give $\Phi_\omega(X)\ge\Phi_\omega(Y)$, which is false. If $X\triangle Y$ is null, $X\preceq^\ast_\omega Y$ by (L1).
\end{proof}

\begin{lemma}\label{lem:uniform}
If (L3) holds for $(X,Y)$, there is $\delta>0$ such that (L3) holds for $(X,Y')$ for every $Y'\in\M$ with $d(Y',Y)<\delta$.
\end{lemma}

\begin{proof}
Let $Z$ satisfy (L3) for $(X,Y)$, $\gamma\coloneqq\gamma(X,Z)>0$ and $\eta\coloneqq\frac14\gamma-\nu(P(Y,Z))>0$. Let $Y'\in\M$ with $d(Y',Y)<\delta$, where $\delta$ is fixed below. Put $\Sigma\coloneqq P(X,Z)\setminus P(X,Y')$, a symmetric set. If $g\in\Sigma$, then $g\in P(X,Y)$, so $g$ or $-g$ lies in $X\triangle Y$, and neither lies in $X\triangle Y'$; hence $g$ or $-g$ lies in $Y\triangle Y'$. Thus $\Sigma\subseteq P(Y,Y')$, and by Lemma~\ref{lem:separable}
\[
\mu(\Sigma)\le2\delta,\qquad \nu(\Sigma)\le4\delta.\tag{1}
\]
We construct $Z'$ with $J(X,Z')$, $P(X,Z')\subseteq P(X,Z)\setminus\Sigma$, and $\nu(P(Z,Z'))\le\kappa\delta$, where $\kappa$ depends only on $Z$.

\emph{Case $J_\GU(X,Z)$.} Let $Z'\coloneqq(Z\setminus\Sigma)\cup(X\cap\Sigma)$. Then $Z'$ agrees with $Z$ off $\Sigma$ and with $X$ on $\Sigma$, so $P(X,Z')=P(X,Z)\setminus\Sigma$ and $P(Z,Z')\subseteq\Sigma$; $\nu(P(Z,Z'))\le4\delta$. For every $\omega$ and almost every $g\in P(X,Z')$ with $g(\omega)\ne0$, $\rho_\omega(Z',g)=\rho_\omega(Z,g)$, so the transition is required strict. By Lemma~\ref{lem:Phi}(b), $\gamma(X,Z')\ge\gamma-4\delta$; if $4\delta<\gamma$, then $\gamma(X,Z')>0$, hence $\mu(P(X,Z'))>0$ and $J_\GU(X,Z')$. Put $\kappa\coloneqq4$.

\emph{Case of an averaging, weighted sure-loss withdrawal or concentration instance.} Let $A_1,\dots,A_n$ be its sets ($G_1,G_2,H$; or $G_1,\dots,G_k$; or $G,H$), with measures $m_\ell\coloneqq\mu(A_\ell)$, $m_{\min}\coloneqq\min_\ell m_\ell$, $M_{\mathrm{tot}}\coloneqq\sum_\ell m_\ell$, centres $c_\ell$, and $M_0\coloneqq\max_\ell\|c_\ell\|+r$. Then $P(X,Z)=\bigcup_\ell(A_\ell\cup-A_\ell)$. Let $2\delta<m_{\min}$ and $\theta\coloneqq1-2\delta/m_{\min}\in(0,1)$. Since $\mu(A_\ell\cap\Sigma)\le2\delta$ by (1),
\[
\mu(A_\ell\setminus\Sigma)\ge m_\ell-2\delta\ge m_\ell-2\delta\,\tfrac{m_\ell}{m_{\min}}=\theta m_\ell,
\]
so (M1) gives measurable $A_\ell'\subseteq A_\ell\setminus\Sigma$ with $\mu(A_\ell')=\theta m_\ell$. Replacing each $A_\ell$ by $A'_\ell$ and keeping all centres and $r$ gives an instance of the same type: for averaging and concentration the inequalities of Definitions~\ref{def:AV} and~\ref{def:CONC} involve only centres and $r$, and the measures are $\theta m,\theta m,2\theta m$ and $\theta m,\theta m/t$; for weighted sure-loss withdrawal, $\sum_i\theta m_if_i\le-(\theta\varepsilon+r\sum_i\theta m_i)\one$. Let $Z'$ be the result of this instance applied to $X$. Then $J(X,Z')$, and $P(X,Z')=\bigcup_\ell(A'_\ell\cup-A'_\ell)\subseteq P(X,Z)\setminus\Sigma$ because $\Sigma$ is symmetric. $P(Z,Z')=\bigcup_\ell\bigl((A_\ell\setminus A'_\ell)\cup-(A_\ell\setminus A'_\ell)\bigr)$, on which $\|g\|<M_0$, and $\mu(A_\ell\setminus A'_\ell)=(1-\theta)m_\ell=2\delta m_\ell/m_{\min}$. Hence
\[
\nu(P(Z,Z'))\le2M_0\cdot2\sum_\ell(1-\theta)m_\ell=\frac{8M_0M_{\mathrm{tot}}}{m_{\min}}\,\delta\eqqcolon\kappa\delta.
\]

\emph{Both cases.} By Lemma~\ref{lem:Phi}(b), $\Gamma_\omega(X,Z')\ge\Gamma_\omega(X,Z)-\nu(P(Z,Z'))$ for every $\omega$, so $\gamma(X,Z')\ge\gamma-\kappa\delta$. Since $P(Y',Z')\subseteq P(Y',Y)\cup P(Y,Z)\cup P(Z,Z')$ and by Lemma~\ref{lem:separable},
\[
\nu(P(Y',Z'))\le4\delta+\nu(P(Y,Z))+\kappa\delta.
\]
Also $P(X,Z')\subseteq P(X,Z)\setminus\Sigma\subseteq P(X,Y')$ by the definition of $\Sigma$. Choose $\delta>0$ with $\kappa\delta<\gamma$, $2\delta<m_{\min}$ in the second case, and $4\delta+\frac54\kappa\delta<\eta$. Then
\[
\nu(P(Y',Z'))\le\tfrac14\gamma-\eta+4\delta+\kappa\delta<\tfrac14(\gamma-\kappa\delta)\le\tfrac14\gamma(X,Z'),
\]
so $Z'$ satisfies (L3) for $(X,Y')$.
\end{proof}

\begin{theorem}\label{thm:R-in-class}
$L^\ast\in\K_t$.
\end{theorem}

\begin{proof}
\emph{Axiom G.} (i) Let the hypothesis hold at $\omega$. If $\mu(P(X,Y))=0$ then (L1). Otherwise, for almost every $g\in P(X,Y)$ with $g(\omega)\ne0$ the responses differ (Definition~\ref{def:changed}), and a required weak transition $(b,a)$ has $a\rhd b$; so (L2) holds. (ii) Then (L2) holds and $X\triangle Y$ is not null, so $Y\prec^\ast_\omega X$ by Lemma~\ref{lem:monotone}.

\emph{Axioms AV, WSL, CONC.} For an instance $(X,Z)$, $J(X,Z)$ holds, and $Z$ itself satisfies (L3) for $(X,Z)$, since $\nu(P(Z,Z))=0<\frac14\gamma(X,Z)$ by Lemma~\ref{lem:Phi}(e). So $Z\preceq^\ast_\omega X$ for every $\omega$; $X\triangle Z\supseteq G_1$ (or $G$) is not null, so $Z\prec^\ast_\omega X$ by Lemma~\ref{lem:monotone}.

\emph{Axiom I.} Let $X,Y,X',Y'$ and $P$ be as in Axiom I. It suffices to show $Y\preceq^\ast_\omega X\Rightarrow Y'\preceq^\ast_\omega X'$; the converse follows by exchanging the primed and unprimed models, and the second equivalence by exchanging the roles of $X$ and $Y$ and of $X'$ and $Y'$, since the hypotheses are symmetric in these roles. (L1): $X\triangle Y\subseteq P$ is null iff $P$ is null iff $X'\triangle Y'$ is null. (L2): the condition refers only to $P$ and to the responses of $X$ and $Y$ to classes in $P$, which coincide with those of $X'$ and $Y'$ almost everywhere on $P$. (L3): let $Z$ satisfy (L3) for $(X,Y)$, put $Q\coloneqq P(X,Z)\subseteq P$, and let $Z'\coloneqq(X'\setminus Q)\cup(Z\cap Q)$. Then $Z'$ agrees with $X'$ off $Q$ and with $Z$ on $Q$, and $X'$ agrees with $X$ on $Q$; hence $P(X',Z')=Q$ and $X',Z'$ give the same responses on $Q$ as $X,Z$. Every condition defining $J$ refers only to $Q$, to these responses, and to subsets of $Q$ lying in $U(\cdot)$ or $S(\cdot)$ of the first model; so $J(X',Z')$, and $Z'$ is obtained from $X'$ by the same sets as $Z$ from $X$. By Lemma~\ref{lem:Phi}(a) with $A=Q$, $\Gamma_\omega(X',Z')=\Gamma_\omega(X,Z)$ for every $\omega$, so $\gamma(X',Z')=\gamma(X,Z)$. Off $P$, $Y=X=Z$ and $Y'=X'=Z'$; on $P$, $Y'$ agrees with $Y$ and $Z'$ with $Z$; so $P(Y',Z')=P(Y,Z)$. Finally $P(X',Z')=Q\subseteq P=P(X',Y')$. Hence $Z'$ satisfies (L3) for $(X',Y')$.

\emph{Axiom UR.} Let $J(X,Z)$. Then $Z$ satisfies (L3) for $(X,Z)$. By Lemma~\ref{lem:uniform} there is $\delta>0$ such that (L3) holds for $(X,Y)$ whenever $d(Y,Z)<\delta$. For such $Y$, $Y\preceq^\ast_\omega X$ for every $\omega$, and $\Phi_\omega(Y)>\Phi_\omega(X)$ by Lemma~\ref{lem:monotone}, so $X\triangle Y$ is not null and $Y\prec^\ast_\omega X$ for every $\omega$.
\end{proof}

\begin{theorem}\label{thm:protection}
Every $T\in\C$ is undominated under $L^\ast$.
\end{theorem}

\begin{proof}
Suppose $Y$ dominates $T$ under $L^\ast$. Then $T\triangle Y$ is not null, since otherwise $T\preceq^\ast_\omega Y$ for every $\omega$ by (L1). By Lemma~\ref{lem:noJ}, (L3) fails for $(T,Y)$. Hence (L2) holds for $(T,Y)$ at every $\omega\in\Omega$. As in the proof of Lemma~\ref{lem:noJ}, $K(T)$ is null, so $\rho_\omega(T,g)\ne\con$ for almost every $g$. A pair $(b,a)$ with $a\rhd b$ and $b\ne\con$ is a required strict transition. Hence, for every $\omega$ and almost every $g\in P(T,Y)$ with $g(\omega)\ne0$, $(\rho_\omega(T,g),\rho_\omega(Y,g))$ is a required strict transition; and $\mu(P(T,Y))>0$. So $J_\GU(T,Y)$ holds, contradicting Lemma~\ref{lem:noJ}.
\end{proof}

\section{The main theorem}\label{sec:main}

\begin{theorem}\label{thm:main}
Assume (M1)--(M4) and fix $t>1$.
\begin{enumerate}[label=(\alph*)]
\item $L^\ast\in\K_t$, and every almost coherent model is undominated under $L^\ast$.
\item For every $L\in\K_t$, every $X\in\M\setminus\C$ is dominated under $L$ by a model $Z$ with $J(X,Z)$, and $Z\prec_\omega X$ for every $\omega$.
\item Every $L\in\K_t$ satisfies countable selection on $\M\setminus\C$, with models $c^j_p$ satisfying (F3).
\item For every family $\Lf\subseteq\K_t$ with $L^\ast\in\Lf$, conditions (E), (A), (S) and (T) hold. Consequently the rationally permissible models relative to $\Lf$ are exactly the almost coherent models, and the irrational models are exactly the models that are not almost coherent.
\end{enumerate}
\end{theorem}

\begin{proof}
(a) Theorems~\ref{thm:R-in-class} and~\ref{thm:protection}. (b) Corollary~\ref{cor:direct}. (c) Theorem~\ref{thm:selection}. (d) (E) holds because every ranking is defined on all of $\M$. (A): $\C_L\subseteq\C$ for every $L$, and $\C_{L^\ast}=\C$ by (a). (S) and (T) hold by (c). Theorem~\ref{thm:ldg} with $\C^\ast\coloneqq\C$ gives the conclusion.
\end{proof}

\section{Models without concentration}\label{sec:convex}

\begin{definition}\label{def:convex-coherent}
A model $X$ is \emph{convex-coherent} if there is a closed convex $R\subseteq E$ with $E_+\subseteq R$, $R+E_+\subseteq R$, $R\cap E_{--}=\varnothing$ and $X\triangle\interior R$ null. $\C^0$ is the set of convex-coherent models. $J^0(X,Z)$ holds if $J_\GU(X,Z)$ holds or $(X,Z)$ is an averaging or a weighted sure-loss withdrawal instance. $\K^0$ is the class of rankings satisfying Axioms G, AV, WSL, I, and Axiom UR with $J^0$ in place of $J$.
\end{definition}

\begin{theorem}\label{thm:convex}
Assume (M1)--(M4).
\begin{enumerate}[label=(\alph*)]
\item If there is no $Z$ with $J^0(X,Z)$, then $X\in\C^0$.
\item If $T\in\C^0$, there is no $Z$ with $J^0(T,Z)$.
\item Let $L^0$ be defined as in Definition~\ref{def:R} with $J^0$ in place of $J$. Then $L^0\in\K^0$, and every $T\in\C^0$ is undominated under $L^0$.
\item For every family $\Lf\subseteq\K^0$ with $L^0\in\Lf$, conditions (E), (A), (S) and (T) hold with $\C^0$ in place of $\C$, and the rationally permissible models are exactly the convex-coherent models.
\end{enumerate}
\end{theorem}

\begin{proof}
(a) Steps 0, 1, 4 and 5 of the proof of Theorem~\ref{thm:classification} use only $J_\GU$, averaging and weighted sure-loss withdrawal instances; Step 4 uses (M4) for the set $R$, which by Step 0(d) and Consequences (i), (ii) of Step 1 is closed, convex, satisfies $E_+\subseteq R$, $R+E_+\subseteq R$ and $R\cap E_{--}=\varnothing$. Step 5 uses only $0\notin\interior R$, which follows from $R\cap E_{--}=\varnothing$. They give $X=U(X)$ and $U(X)\triangle\interior R$ null, so $X\in\C^0$.

(b) Let $R$ be as in Definition~\ref{def:convex-coherent} and $T\triangle\interior R$ null. $\interior R$ contains no $g\le0$ (as in Step 6, (C2)), and $g,-g\in\interior R$ would give $0=\frac12g+\frac12(-g)\in\interior R$; so $K(T)$ is null and $U(T)=\interior R$ up to null sets. \emph{Averaging}: choose $f_i'\in G_i\cap\interior R$; for $h'\in H$, $h'\ge\frac12(f_1'+f_2')+\varepsilon\one$, and $\frac12(f'_1+f'_2)\in\interior R$, so $h'\in\interior R$ as in Step 6, (C4); hence $H\subseteq\interior R$ and $\mu(H)=0$, a contradiction. \emph{Weighted sure-loss withdrawal}: with $f'_i\in G_i\cap\interior R$ and $M\coloneqq\sum_im_i$, $a\coloneqq\sum_i\frac{m_i}{M}f_i'\in\interior R$ by convexity, and $a\le-\frac{\varepsilon}{M}\one$, contradicting $R\cap E_{--}=\varnothing$. \emph{$J_\GU$}: $E_+$ is closed and $\interior E_+=E_{++}$, since $B(g,\inf g)\subseteq E_+$ for $g\in E_{++}$, and $g-\tau\one\notin E_+$ for every $\tau>0$ if $g\in E_+$ and $\inf g=0$. By (M4) applied to $E_+$, $\bd E_+=E_+\setminus E_{++}$ is null, and $E_{++}\subseteq\interior R$. Proceed as in the proof of Lemma~\ref{lem:noJ}, with $\interior R$ in place of $D$, adding $\bd E_+\cup-\bd E_+$ to $N$. If $\interior R\cap[g]=\{g'\}$, then $g'\not\le0$, so $g'(w)>0$ for some $w\in W$. If $\interior R\cap[g]=\varnothing$, then $g\notin E_{++}$ and $-g\notin E_{++}$, and $g,-g\notin\bd E_+$, so neither $g$ nor $-g$ lies in $E_+$: $g$ takes a strictly positive and a strictly negative value. The rest of that proof applies verbatim.

(c) The proofs of Lemmas~\ref{lem:Phi}, \ref{lem:monotone}, \ref{lem:uniform} and Theorem~\ref{thm:R-in-class} use only the clauses of $J$ that define $J^0$ when $J^0$ replaces $J$, and so hold for $L^0$. The proof of Theorem~\ref{thm:protection} with (b) in place of Lemma~\ref{lem:noJ}, and with $K(T)$ null by (b), gives the protection.

(d) (S) and (T) with $\C^0$ in place of $\C$: Lemma~\ref{lem:general-selection} with $(\M',d')\coloneqq(\M,d)$, $\C'\coloneqq\C^0$ and $J'\coloneqq J^0$, whose hypotheses are (a) and Axiom UR for $J^0$. (A) with $\C^0$: $\C^0\cap\U_{L^0}=\C^0$ by (c). (E) holds since every ranking is defined on $\M$. Theorem~\ref{thm:ldg} with $\C^\ast\coloneqq\C^0$ gives the conclusion.
\end{proof}

\begin{thebibliography}{9}
\bibitem{LDG} J.~Konek, \emph{Dominance Generalised} (countable selection version), manuscript, 30 October 2026. Numbering as in that version.
\end{thebibliography}

\end{document}
